% chercher des documents LaTeX dans styles, corps et bib
\makeatletter\def\input@path{{styles/}{corps/}}\makeatother

% Utiliser le style rapport.cls
\documentclass{rapport}
\usepackage{fancyhdr}

\fancyhead[L,C]{}
\fancyhead[R]{Rapport final de Projet Informatique}
\fancyfoot[C]{Alexandre Froehlich \& Guillaume Leinen}
\fancyfoot[L]{U.E. 2.4}
\fancyfoot[R]{\thepage}
\renewcommand{\headrulewidth}{0.4pt}
\renewcommand{\footrulewidth}{0.4pt}


\title{Rapport final de Projet Informatique\\\og La Bataille Brestoise\fg}

\author{Alexandre \textsc{Froehlich} \& Guillaume \textsc{Leinen}\\
  \texttt{alexandre.froehlich@ensta-bretagne.org}
  \texttt{guillaume.leinen@ensta-bretagne.org}\\
	Groupe 1}
\date{\today}
\doctype{Projet Info.\\ U.E. 2.4.}
\promo{FISE 2022}
\etablissement{\textsc{Ensta} Bretagne\\2, rue Fran�ois Verny\\
  29806 \textsc{Brest} cedex\\\textsc{France}\\Tel +33 (0)2 98 34 88 00\\ \url{www.ensta-bretagne.fr}}
\logoEcole{\includegraphics[height=4.2cm]{logo_ENSTA_Bretagne_Vertical_CMJN}}

% M�tadonn�es � include dans le pdf
\hypersetup{%
  pdftitle={Rapport premi�re phase},
  pdfauthor={Alexandre Froehlich, Guillaume Leinen},
  pdfkeywords={},
  bookmarksnumbered,
  pdfstartview={FitH},
  citecolor=blue,
  breaklinks=true
}%

% creer le glossaire
\makeglossaries
% creer l'index
\makeindex

\pagestyle{fancy}

\begin{document}
% D�but du pr�ambule
\frontmatter
% inclure la liste des acronymes
\newglossaryentry{iot}{type=glo,
  text={\textsc{IoT}},
  name={Internet of Things},
  description={Ensembles des objets dits "connect�s". Cela va du smartphone, � la t�l�vision, au frigo...}}

\newglossaryentry{ia}{type=glo,
	text={\textsc{IA}},
	name={Intelligence Artificielle},
	description={Algorithme ou machine qui a pour but de simuler une intelligence}}

\newglossaryentry{enigma}{type=glo,
	text={\textsc{Enigma}},
	name={Enigma},
	description={Machine servant au chiffrement et d�chiffrement de messages, d'informations principalement employ�e par les Allemands}}

\newglossaryentry{alphago}{type=glo,
	text={\textsc{AlphaGo}},
	name={AlphaGo},
	description={Programme informatique qui a pour but de jouer au Go, il a remport� de nombreux prix face � de grands joueurs humains}}

\newglossaryentry{resneurone}{type=glo,
	text={\textsc{R�seau Neuronal}},
	name={R�seau Neuronal},
	description={Programme ou syst�me que l'on va entra�ner � reconnaitre des motifs, des tendances, ...}}

\newglossaryentry{dataset}{type=glo,
	text={\textsc{dataset}},
	name={Dataset},
	description={Ensemble de donn�es. Peut �tre des images, des vid�os, du texte, ...}}

\newglossaryentry{pixel}{type=glo,
	text={\textsc{pixel}},
	name={Pixel},
	description={Unit� �l�mentaire composant une image, poss�de une unique couleur}}

\newglossaryentry{MNIST}{type=glo,
	text={\textsc{MNIST}},
	name={MNIST Dataset},
	description={Ensemble de donn�es du Mixed National Institute of Standards and Technology, un dataset contenant des chiffres �crits � la main}}

\newglossaryentry{CNN}{type=glo,
	text={\textsc{CNN}},
	name={Convolutional Neural Network},
	description={R�seau neuronal dans lequel le motif de connexion entre les neurones est inspir� du cortex visuel des animaux}}

\newglossaryentry{RCNN}{type=glo,
	text={\textsc{R-CNN}},
	name={Regions Convolutional Neural Network},
	description={Un CNN qui ne travaille que sur une r�gion donn�e d'une information en entr�e}}

\newglossaryentry{TPU}{type=glo,
	text={\textsc{TPU}},
	name={Tensor Processing Unit},
	description={Mat�riel sp�cifiquement con�u par Google pour calculer les tenseurs, structure math�matique au c?ur de ses algorithmes}}

\newglossaryentry{framework}{type=glo,
	text={\textsc{framework}},
	name={Framework},
	description={Cadre de travail, c'est souvent un ensemble de m�thodes et de fonctions d�finies dans une biblioth�que}}

\newglossaryentry{npl}{type=glo,
	text={\textsc{Natural Processing Language}},
	name={Natural Processing Language},
	description={Ensemble de techniques qui ont pour but d'harmoniser le langage humain et le langage machine}}

\newglossaryentry{appsuper}{type=glo,
	text={\textsc{apprentissage supervis�}},
	name={Apprentissage supervis�},
	description={Apprentissage de pr�diction avec un dataset d'entr�e et le r�sultat suppos� en sortie}}

\newglossaryentry{appnonsuper}{type=glo,
	text={\textsc{apprentissage non-supervis�}},
	name={Apprentissage non-supervis�},
	description={Apprentissage sans pr�diction, on laisse le r�seau neuronal "explorer"}}

\newglossaryentry{apprenfor}{type=glo,
	text={\textsc{apprentissage par renforcement}},
	name={Apprentissage par renforcement},
	description={Apprentissage par s�lection ou r�compense. On r�compense le r�seau neuronal pour les r�ponses les plus adapt�es qu'il donne}}

\newglossaryentry{deep}{type=glo,
	text={\textsc{deep learning}},
	name={Deep Learning},
	description={Apprentissage profond. Devenu un terme g�n�rique pour parler d'un algorithme poss�dant un r�seau neuronal}}

\newglossaryentry{transfer}{type=glo,
	text={\textsc{transfer learning}},
	name={Transfer Learning},
	description={Concept qui prend un r�seau neuronal pr�-entra�n� et qui l'exploite sur une application similaire sans passer par la phase d'apprentissage exclusif}}

%%% Local Variables: 
%%% mode: latex
%%% TeX-master: "../guide"
%%% End: 

% creer le titre ici
\maketitle

\clearpage

% table des matieres
\tableofcontents

\clearpage

% Partie principale
\mainmatter
% chargement du fichier corps.tex

\section{Description g�n�rale du probl�me}
\subsection{Introduction}

Pour ce projet d'informatique, nous avons d�cid� de prendre un sujet autre que ceux propos�s par nos enseignants.  Ayant fait tous les deux une ann�e z�ro dans la marine nationale, un sujet proche de ce que l'on avait v�cut semblait la solution �vidente.

Nous sommes donc parti sur un jeu de type bataille navale, en temps r�el et multijoueurs que nous avons nomm� :

\begin{figure}[!h]
	\centering
	\includegraphics[width=0.5\linewidth]{../imgs/logo}
	\caption{Logo du projet d'informatique}
	\label{fig:logo}
\end{figure}

Pour travailler en collaboratif sur le code, nous avons utilis� un repository sur github, le code est accessible au lien suivant :\\ \url{https://github.com/NightlySide/La-Bataille-Brestoise}.

\subsection{Les diff�rentes pistes envisag�es}

Pour le choix de l'interface Homme-Machine (IHM) le choix s'est rapidement port� sur \textbf{PyQt5} pour sa polyvalence et sa programmation �v�nementielle. Concernant la connexion entre le serveur et les clients,  nous sommes tout d'abord partis sur la biblioth�que \textbf{sockets} qui permet de cr�er un tunnel simple entre deux clients. Nous avons finalement choisi \textbf{asyncio}, impl�ment�e depuis python3, qui permet de g�rer en asynchrone toutes les requ�tes n�cessaires au projet, et ce avec plusieurs clients en simultan�.

\subsection{Diff�rentes hypoth�ses r�ductrices choisies}

Pour r�duire la taille du projet et surtout le rendre r�alisable, nous avons choisi d'�tablir quelques r�ductions pour cette premi�re partie : 
\begin{itemize}
	\item Le positionnement et le d�placement des joueurs et entit�s se fait une une grille de $8\times 8$ pixels.
	\item Les joueurs ne sont pas visibles par les autres joueurs.
	\item Le chat ne peut accueillir que du texte, donc pas de commande\\ (du type \emph{/help}.
	\item La vitesse de rafraichissement, ou le nombre de boucles de jeu par seconde est ind�pendant selon la classe.
	
\end{itemize}

\clearpage
\section{Application proprement dite}
\subsection{Pr�sentation g�n�rale du programme}
Le programme est scind� en deux partie : le client et le serveur avec chacun son script pour le lancer. Le script \emph{test$\_$suite.py} permet de lancer tous les tests unitaires contenus dans le dossier \emph{tests}.

\subsubsection{Le serveur}
Lancer le serveur avec le script \emph{serveur.py}, cr�e un serveur sur l'h�te \emph{localhost} (127.0.0.1 ou 0.0.0.0). sur le port 25566.
Le serveur va ensuite accueillir les nouveaux joueurs, organiser le chat et g�rer toutes les entit�s du jeu. Une fois lanc� et op�rationnel le terminal de commande ressemble � cela :

\begin{figure}[H]
	\centering
	\includegraphics[width=0.8\linewidth]{../imgs/serveur_screen}
	\caption{Ecran de fonctionnement du serveur}
	\label{fig:serveur_screen}
\end{figure}

Une fois d�marr� le serveur affichera les informations utiles telles que le nom de la machine sur laquelle il est ex�cut�, son ip locale, le port de connexion et le nombre de joueurs maximum (pas encore impl�ment�).

\subsubsection{Le client}
Une fois le serveur lanc� en local, on peut soit se connecter � ce dernier, soit se connecter � un serveur que nous laisserons ouvert chez nous, pr�-enregistr� dans la liste des serveurs disponibles.

En d�marrant le client avec le script \emph{client.py}, on se trouve devant l'�cran de connexion du projet. (cf. \ref{fig:ecran_connexion})

Sur celui-ci l'utilisateur pourra retrouver les serveurs pr�-enregistr�s ainsi qu'une boite pour indiquer son pseudonyme, un bouton pour se connecter au serveur s�lectionn�. 

Enfin l'utilisateur trouvera  un bouton "Serveur local" pour entrer manuellement les informations de connexion au serveur (cf. \ref{fig:serveur_local})

\begin{figure}[H]
	\centering
	\includegraphics[width=0.8\linewidth]{../imgs/connexion_screen}
	\caption{Ecran de connexion client}
	\label{fig:ecran_connexion}
\end{figure}

\begin{figure}[H]
	\centering
	\includegraphics[width=0.5\linewidth]{../imgs/serveur local}
	\caption{Boite de dialogue de connexion manuelle � un serveur}
	\label{fig:serveur_local}
\end{figure}

\clearpage

Une fois connect� au serveur, l'utilisateur est face � l'�cran de jeu (cf. \ref{fig:ecran_jeu}) qui comporte une surface de jeu, que nous appellerons \emph{canvas}, sur laquelle sera dessin�e la carte, les entit�s et les joueurs. On y retrouve un radar permettant de d�tecter les entit�s aux alentours (attention pas enti�rement fonctionnel pour l'instant), et une chatbox pour communiquer avec les autres joueurs pr�sent sur le serveur.

\begin{figure}[H]
	\centering
	\includegraphics[width=1\linewidth]{../imgs/game_screen}
	\caption{Ecran de jeu du projet}
	\label{fig:ecran_jeu}
\end{figure}

Pour cette premi�re phase, nous nous sommes concentr� sur la structure du jeu. Par cons�quent les contr�les joueur sont limit�s. Pour se d�placer l'utilisateur pourra utiliser les fl�ches du clavier d�s que le focus est sur le canvas (un click souris sur le canvas suffit � r�cup�rer le focus).

\subsection{Description des principales classes/m�thodes}

Pour une meilleure lisibilit�, vous pouvez retrouver les diagrammes de classes en annexe (cf. \ref{fig:annexe_diagramme_complet})

% d�crire uniquement les m�canismes fondamentaux
\subsubsection{Classes de jeu}

En premier lieu, il convient de d�crire les classes directement li�s au fonctionnement du jeu.  (cf .\ref{fig:annexe_diagramme_client_jeu})

La classe "Entit�" est la base d'une joueur, qu'il soit r�el ou contr�l� par la machine. Elle va notamment avoir des attributs de points de vie, de position dans la map ou encore d'arme actuellement s�lectionn� par le joueur.

Une entit� peut �tre "joueur" ou "pnj". Une entit� interagit avec les autres entit�s par le biais du client de jeu, les entit� "pnj" sont elles contr�l�s par le serveur. Une entit� � un instant t contr�le un b�timent et s�lectionne une arme disponible sur ce b�timent. La classe "Batiment" contient toutes les statistiques et donn�es sur un b�timent, cela peut �tre son niveau, les points de vie maximaux, ou encore la port�e de d�tection. La classe "Armes" quand � elle d�finie les donn�es des diff�rentes armes disponibles.  (cf. \ref{fig:annexe_diagramme_mecanismes})

\subsubsection{Classe de l'IHM}


Dans l'IHM, nous avons int�gr� plusieurs �l�ments qui n�cessite leur propre classe. 

La classe "Radar" est une classe h�rit� de Qwidget. Elle g�re le radar de d�tection des autres entit�s ennemies. \\La r�gle du jeu veut que si un ennemi se trouve dans notre port�e de d�tection et que cette ennemi a une dissimulation (impl�ment�e plus tard) sup�rieur � la distance qui nous s�pare d'elle, cette ennemi sera affich� par un point rouge sur le radar. \\Le radar balaye l'espace � 90�/s. On garde les traces d'un ennemie d�tect� en rouge puis en gris lorsqu'elles s'appr�tent � disparaitre du radar. 


\begin{figure}[H]
	\centering
	\includegraphics[width=0.5\linewidth]{../imgs/radar}
	\caption{Aper�u du Radar}
	\label{fig:radar}
\end{figure}


La classe "Chatbox" comme son nom l'indique renvoi � la chat box en bas � droite du canevas de jeu. Elle est h�rit�e de QTextBrowser. Cette classe existe car on doit pouvoir �crire sur le chat pendant que le jeu tourne, il est donc n�cessaire de cr�er un autre thread, sur lequel tourne la chatbox, et dans lequel Qt va piocher les informations lorsque n�cessaire.\\ Elle contient une liste de lignes qui est mise � jour avec les messages des autres joueurs. Un signal Qt permet d'envoyer un message en appuyant sur la touche Entr�e. 

\subsubsection{Classe du client/serveur }

Pour g�rer la connexion asynchrone entre le client et le serveur il est n�cessaire d'avoir deux classes suppl�mentaires g�rant ces diffrents aspects. La classe 'TCPClientProtocol' (cf. \ref{fig:annexe_diagramme_client_jeu}) h�rite du protocole de la bilbioth�que \emph{asyncio} et qui fonctionne en �v�nementiel. La classe 'TCPServer' du c�t� serveur (cf. \ref{fig:annexe_diagramme_serveur}), poss�de un fonctionnement similaire. On retrouve ainsi dans les deux classes une fonction pour prendre en charge l'ouverture du tunnel entre le client et le serveur, une fonction pour prendre en charge la d�connexion et enfin deux fonctions pour l'envoi et la r�ception de donn�es.

\section{Conclusion}

\subsection{Diff�rents tests effectu�s}

Pour les tests unitaires, le script \emph{test\_suite.py} permet d'effectuer les tests provenant des classes de \emph{unittest} pr�sentes dans le dossier 'tests'. Ces classes de test vont tester les �l�ments de l'IHM ainsi que certaines m�caniques du jeu.

Pour la connexion client/serveur, nous avons tout d'abord essay� la biblioth�que \emph{sockets} afin de cr�er la connexion entre le client et le serveur. Cependant cette derni�re ne g�re que difficilement les connexions multiples du c�t� du serveur. Il est vrai que l'on aurait pu utiliser comme avec python 2.7 le module \emph{selectors}, cependant cette m�thode, comme \emph{asyncore} que nous avons aussi essay�, semble dispara�tre avec les versions apr�s 3.6 au profit de \emph{asyncio},.

Pour suivre l'�volution des tests et � fins de d�bug, nous avons impl�ment� un \emph{Logger} dans notre projet. Il va g�rer l'affichage sur le terminal de commande en fonction de la s�v�rit� du message envoy� par le programme, il peut de plus, enregistrer ces logs dans un fichier \emph{.log}.

\subsection{Limitations}

Nous avons rencontr� quelques limitations durant ce projet.

En premier lieu, l'usage d'une mode multijoueur implique une utilisation extensive de la bande passante, ce qui n�cessite d'optimiser les requ�tes afin de ne pas surcharger un r�seau.

Ensuite le fait que Pickle ne g�re pas directement QT, en raison du fait qu'il ne g�re pas les objets cod�s en C. Hors, nos messages sont s�rialis�s � l'aide de \emph{pickle} et envoy�s sur le r�seau. Il faut donc se restreindre � de l'envoi d'objets de 'bases' reconnus par pickle lors de la communication client/serveur.

Une derni�re limitation est le fait que la biblioth�que \emph{pyQt}, n'accepte pas que l'on interragisse avec ses objets � partir d'un thread diff�rent. La solution fut de passer par des classes annexes comme "ChatBox" afin de contourner cette limitation, en ayant d'un c�t� le thread TCP populant le contenu de la ChatBox, et de l'autre pyQt qui r�cup�re ces informations pour l'afficher sur l'IHM (comme expliqu� pr�c�demment et cf. \ref{fig:chatbox})

\begin{figure}[H]
	\centering
	\includegraphics[width=0.7\linewidth]{../imgs/chatbox}
	\caption{Principe de la classe interm�diaire entre deux Threads}
	\label{fig:chatbox}
\end{figure}

Nous avons pu contourner ces limitations et les seuls r�els limitations de notre projet sont actuellement plus de l'ordre de la performance que du fonctionnement. 


\subsection{Perspectives d'am�lioration}

On pourrait am�liorer la dynamique des b�timents, en leur donnant des caract�ristiques d'inertie ou encore de rayon de giration, ce qui aurait �t� plus r�aliste (on ne manoeuvre pas un porte-avions comme un BIN). 
\\On pourrait encore impl�menter un concept de d�g�t critique, avec des �quipements importants du b�timent qui pourrait �tre mis temporairement hors d'usage apr�s un coup bien plac� (armement, radar, propulsion).  \\ Pour am�liorer le temps de calcul on pourra par exemple r�fl�chir � un moyen de passer en calcul vectoriel si possible.

Enfin on pourra impl�menter des sauvegarde de score ou tout simplement des sauvegarde de param�tres utilisateurs dans des fichiers internes au projet, probablement en JSON ou en YAML.




\clearpage
\appendix
\pagestyle{empty}
% \section{A ajouter s'il le faut}


%\clearpage
%\phantomsection
%\label{sec:index}
%\addcontentsline{toc}{section}{Index}
%\printindex

%\printglossaries


%\clearpage

\end{document}

%%% Local Variables: 
%%% mode: latex
%%% TeX-master: t
%%% End: 
